\documentclass[11pt,a4paper]{article}

\usepackage[T1]{fontenc}
\usepackage[utf8]{inputenc}
\usepackage[margin=2cm]{geometry}
\usepackage{microtype}
\usepackage{amsmath,amssymb,amsthm,mathtools}
\usepackage{bm}
\usepackage{physics}
\usepackage{xcolor}
\usepackage{markdown}
\usepackage{hyperref}
\hypersetup{
  colorlinks=true,
  allcolors=blue,
}

\newcommand{\lean}{\texttt{Lean}}
\newcommand{\mathlib}{\textsf{Mathlib}}

% \title{\lean-verified Quantum Information Theory\\QIC891}
% \author{RRS}
% \date{Fall 2026}

\begin{document}

% \maketitle

\noindent
\LARGE\textbf{\href{https://rodolfor-s.github.io/teaching/}{\lean-verified Quantum Information Theory (QIC891)}}\\
\large\textbf{Institute for Quantum Computing Fall 2026 course @ Perimeter Institute for Theoretical Physics}

\normalsize
\section*{Course Overview}

\paragraph{Instructor:} Rodolfo R. Soldati. \textbf{Coordinator:} Debbie Leung.

This syllabus presents the course structure and a roadmap for the full lecture
sequence. The course introduces the principles of programming for formal
mathematics, and how to do quantum information theory in \lean.

From the \textbf{\emph{15th of September}} to the \textbf{\emph{1st of
October}}, 2026.

Every \textbf{\emph{Tuesday}} and \textbf{\emph{Thursday}},
\textbf{\emph{10:30--11:50 am}}.

\textbf{\emph{Space Room}} (4th floor), \textbf{\emph{Perimeter Institute for
Theoretical Physics}}.

\small
\begin{description}
    \item[1st Week]
    \item[Lecture 1, 15 Sep:] Motivation, formal verification and proofs, \lean\ and its foundations.
    %
    \begin{itemize}
        \item What is formal verification and why does it matter for mathematics or physics
        \item Setting up \lean4 and VSCode; first proofs (\verb|#check|, \verb|#eval|, basic tactics)
        \item Propositions, types, and the Curry--Howard correspondence
    \end{itemize}
    %
    \item[Lecture 2, 17 Sep:] Core proof writing, syntax, and tactics in \lean4.
    %
    \begin{itemize}
        \item Propositional and predicate logic in \lean\ (\verb|intro|, \verb|apply|, \verb|exact|, \verb|rw|, \verb|simp|)
        \item Working with natural numbers, integers, and basic algebraic structures
        \item Introduction to Mathlib: navigating the library, finding lemmas
        \item Lab: Proving elementary lemmas about sets and functions
    \end{itemize}
    %
    \item[2nd Week]
    \item[Lecture 3, 22 Sep:] Matrix algebra in \mathlib4.
    %
    \begin{itemize}
        \item Vector spaces and linear maps in Mathlib (\verb|LinearMap|, \verb|Module|, \verb|Subspace|)
        \item Inner product spaces and Hilbert spaces (\verb|InnerProductSpace|, \verb|ComplexHilbertSpace|)
        \item Matrices, adjoints, and unitarity
        \item Lab: Proving basic properties of linear operators
    \end{itemize}
    %
    \item[Lecture 4, 24 Sep:] Quantum mechanics and information theory, \texttt{Physlib/Lean-QuantumInfo}.
    %
    \begin{itemize}
        \item Encoding qubits and quantum states as vectors in \verb|C^n|
        \item Formalizing Hermitian and unitary operators
        \item Tensor products of Hilbert spaces (\verb|TensorProduct|)
        \item Lab: Defining the Bell states and verifying their properties
    \end{itemize}
    %
    \item[3rd Week]
    \item[Lecture 5, 29 Sep:] More quantum information, Generalized Quantum Stein's Lemma.
    %
    \begin{itemize}
        \item Density matrices as positive semidefinite operators with unit trace
        \item Completely positive trace-preserving (CPTP) maps
        \item Partial trace and entanglement formalization
        \item Lab: Encoding a simple quantum channel and verifying CPTP conditions
    \end{itemize}
    %
    \item[Lecture 6, 01 Oct:] Outlook.
    %
    \begin{itemize}
        \item Discussion: current state of quantum information formalization in Mathlib and open problems
        \item Resources for going further (e.g., Physlib efforts, research frontiers)
    \end{itemize}
    %
    \item[Lecture 7, 01 Oct:] (Tentative) Student project presentations.
\end{description}

\normalsize
\subsection*{Course Goals}

\begin{itemize}
  \item To be able to set up a \lean\ repository from scratch.
  \item To understand the role of logical verification and its significance in today's mathematics-based research landscape.
  \item To be able to reason about mathematics and derived topics using code.
  \item To be able to check the validity of AI-written \lean\ code.
\end{itemize}

\subsection*{Assessment and Deliverables}

Assessment based on homework exercises, proof development assignments, and a final project that ties together \lean\ and quantum information theory.

\section*{More resources}

\begin{itemize}
    \item \href{https://www.math.uwaterloo.ca/~wcleung/qic891-f2026.html}{Coordinator's (Debbie) website for the parent course.}
    \item \href{https://rodolfor-s.github.io/teaching/}{Instructor's (me) website for \lean-verified Quantum Information Theory.} See this for the most up-to-date information.
    \item \href{https://github.com/rodolfor-s/lean-quinfo-lectures}{Background repository for the course: \texttt{rodolfor-s/lean-quinfo-lectures}.}
    \item \href{}{\lean4}
    \item \href{}{\mathlib4}
    \item \href{}{physlib}
    \item \href{}{cslib}
\end{itemize}

\end{document}
